\documentclass[10pt,a4paper]{amsart}
\usepackage[margin=19mm]{geometry}
\usepackage{amsmath,amssymb,mathtools}
\usepackage[scr=rsfs,cal=euler]{mathalfa}
\usepackage{xfrac}
\usepackage[unicode,hidelinks,bookmarks=false]{hyperref}
\newcommand{\ie}{i.e.\ }
\newcommand{\cf}{cf.\ }
\newcommand{\resp}{respectively\ }
\newcommand{\Subsection}{\subsection}
\providecommand{\nobreakdash}{\nobreak\hbox{-}}
\setlength{\emergencystretch}{2em}
\multlinegap0pt
\usepackage{bm}
\usepackage{bbm}
\usepackage{dsfont}
\usepackage{xspace}
\usepackage{cancel}
\usepackage{mathtools}
\usepackage{amsthm}
\makeatletter
\@ifundefined{theorem}{\newtheorem{theorem}{Theorem}}{}
\makeatother
\newcommand{\Z}{\mathbb{Z}}
\title[Endotrivial complexes]{Endotrivial complexes}
\author{Sam K. Miller}
\date{Source: arXiv:2309.12138v4 (CC BY 4.0)}
\begin{document}
\maketitle

\noindent\emph{Source and licence.} Endotrivial complexes. Version arXiv:2309.12138v4 (CC BY 4.0), distributed under the Creative Commons Attribution 4.0 International licence, as recorded in the arXiv licence metadata for this exact version and shown on the abstract page https://arxiv.org/abs/2309.12138v4. Full rights notice: licenses/arxiv-2309.12138-CC-BY-4.0.txt.\par\smallskip

\noindent\textit{Editorial scope.} Counted unit: the Theorem whose source label is \texttt{endotrivdef2} in the article (source0.tex lines 459--461, source label \texttt{endotrivdef2}), delivered together with the article's own complete original proof of it (source0.tex lines 464--478, one \texttt{proof} environment of 15 lines). No mathematics is written, repaired or re-derived: statement and proof are the article's own text line for line. Required context retained uncounted: permsplit (lines 282--284). Every definition, display and label of those lines is retained unchanged. Omitted from this package and not redistributed: 1--281, 285--458, 462--463, 479--1225. Nothing inside a retained excerpt is omitted or altered: every retained body is delivered exactly as the article's lines are written, so no omission receipt of any kind accompanies this package. Citations: the retained excerpts carry no citation command at all, so no bibliography entry of the article is transcribed and no reference is redistributed. Reference adaptation: none was needed and none was made; every reference inside the delivered unit resolves inside it, so no unresolved reference or citation is produced. Printed slips kept literal: no printed word, symbol, hypothesis or number of the counted statement or of its proof was altered. Layout and class: the article's own class is \texttt{amsart} with packages including geometry, amsmath, amsthm, amsfonts, amssymb, amscd, lastpage, fancyhdr, mathrsfs, xcolor, graphicx, listings, hyperref, enumitem; the wrapper is the lane's pinned \texttt{amsart} editorial wrapper at 10pt on A4, which restates the theorem-like environments the retained text needs and adds the article's own macro definitions that the retained text needs (\texttt{\textbackslash Z}), with extra wrapper packages (bm, bbm, dsfont, xspace, cancel, mathtools). The delivered numbering is the unit's own. Assets: no retained span contains a figure environment, a graphics-inclusion command, a PDF-page inclusion, a table or a TikZ picture; no image file is copied into this package, the package declares no asset dependency and no image file is redistributed with it. Licence: version arXiv:2309.12138v4 of this article is distributed under the Creative Commons Attribution 4.0 International licence, as recorded in the arXiv licence metadata for this exact version and shown on the abstract page https://arxiv.org/abs/2309.12138v4. A full rights notice is delivered at licenses/arxiv-2309.12138-CC-BY-4.0.txt. Characters: every retained excerpt is pure ASCII, so the delivered engine, XeLaTeX with its default Latin Modern text font, reports no missing character.
	\begin{theorem}\label{permsplit}
		Let $M, N$ be $p$-permutation $kG$-modules, and let $f: M \to N$ be an injective (resp. surjective) homomorphism. Then $f$ is split injective (resp. surjective) if and only if $f(P)$ is injective (resp. surjective) for all $p$-subgroups $P\in s_p(G)$.
	\end{theorem}

	\begin{theorem}\label{endotrivdef2}
		Let $C \in Ch^b({}_{kG}\textbf{triv})$. Then $C$ is endotrivial if and only if for all $p$-subgroups $P \in s_p(G)$, $C(P)$ has nonzero homology concentrated exactly in one degree, with the nontrivial homology having $k$-dimension one.
	\end{theorem}

	\begin{proof}
		First, suppose $C$ is endotrivial, then $C(P)$ is endotrivial as well. Since $H_i(C)^* \cong H_{-i}(C^*)$ for all $i \in \Z$, an easy argument using the K{\"u}nneth formula shows that $C(P)$ has nonzero homology concentrated in one degree, with the nonzero homology having $k$-dimension one.

		Conversely, suppose for all $p$-subgroups $P \in s_p(G)$, $C(P)$ has homology concentrated in one degree, with the homology having $k$-dimension one. We show $C\otimes_k C^* \simeq k$. By the K{\"u}nneth formula, we have that $C \otimes_k C^* \cong D$, where $D$ is a chain complex satisfying $ H_0(D) \cong k$. Label the differentials of $D$ by $d_n: D_n \to D_{n-1}$. If $C$ has length $n$, then $D$ has length $2n - 1$, with $n-1$ and $-(n-1)$ the greatest and least indices respectively for which $D_i \neq 0$. Moreover, $D(P)$ has nonzero homology concentrated in degree zero for any $p$-subgroup $P\in s_p(G)$. Indeed, if $i \in \Z$ is the unique integer for which $H_i(C(P))\neq 0$, then since $H_i(C(P))^* \cong H_{-i}(C(P)^*) \neq 0$
        \begin{align*}
            H_0((C\otimes_k C^*)(P)) &\cong H_0(C(P) \otimes_k C^*(P))\\
            &\cong H_0(C(P) \otimes_k C(P)^*) \\
            &\cong H_i(C(P)) \otimes_k H_{-i}(C(P)^*) \\
            &\cong H_i(C(P)) \otimes_k H_i(C(P))^* \cong  k
        \end{align*}

        It follows similarly by the K\"unneth theorem that $H_i((C\otimes_k C^*)(P)) = 0$ for $i \neq 0.$

		$d_{n-1}(P)$ and $d_{-n + 2}(P)$ are injective and surjective respectively over all $p$-subgroups $P\in s_p(G)$, and therefore by Theorem \ref{permsplit} are split injective and surjective respectively. Thus, we have an isomorphism \[D \cong D' \oplus (D_{n-1} \xrightarrow{\cong} D_{n-1}) \oplus (D_{-n + 2} \xrightarrow{\cong} D_{-n + 2}),\] where $D'$ is a chain complex of length $2n - 3$ for which $D'(P)$ has nonzero homology only in degree zero for all $P \in s_p(G)$. Here, $D_{n-1} \xrightarrow{\cong} D_{n-1}$ is concentrated in degrees $n-1$ and $n-2$, and $D_{-n + 2} \xrightarrow{\cong} D_{-n + 2}$ is concentrated in degrees $-n+2$ and $-n+1$. These two complexes are contractible, so $D \simeq D'$. An induction argument on the length of $C$ yields the homotopy equivalence $D \simeq k$, as desired.
	\end{proof}

\end{document}
